%%%%%%%%%%%%%%%%%%%%%%%%%%%%%%%
\chapter{Dong's high-order outflow boundary condition}
\label{chap:houtflow-bndcond}

\definecolor{nekblue}{HTML}{24527A}
\definecolor{neklightblue}{HTML}{E8F1F7}
\definecolor{nekred}{HTML}{B44B4B}
\definecolor{neklightred}{HTML}{F8EAEA}
\definecolor{nekgreen}{HTML}{3E7654}
\definecolor{neklightgreen}{HTML}{E9F3EC}
\definecolor{nekorange}{HTML}{C87827}
\definecolor{neklightorange}{HTML}{FBF0E3}
\definecolor{nekpurple}{HTML}{76518B}
\definecolor{neklightpurple}{HTML}{F1EAF5}
\definecolor{nekgray}{HTML}{505960}
\definecolor{neklightgray}{HTML}{F1F3F4}
\definecolor{nekdark}{HTML}{263238}

\providecommand{\vect}[1]{\boldsymbol{#1}}
\providecommand{\grad}{\boldsymbol{\nabla}}
\providecommand{\HOutflow}{\textsc{HOutflow}}
\providecommand{\Gammao}{\Gamma_{\!o}}

\tikzset{
  flow/.style={-{Latex[length=2.7mm]},line width=0.75pt,nekgray},
  dataflow/.style={-{Latex[length=2.7mm]},line width=1pt,nekblue},
  stateflow/.style={-{Latex[length=2.7mm]},line width=1pt,nekpurple},
  box/.style={rounded corners=2pt,draw=nekgray,line width=0.65pt,
    fill=white,align=center,inner sep=5pt,font=\small},
  bluebox/.style={box,draw=nekblue,fill=neklightblue},
  redbox/.style={box,draw=nekred,fill=neklightred},
  greenbox/.style={box,draw=nekgreen,fill=neklightgreen},
  orangebox/.style={box,draw=nekorange,fill=neklightorange},
  purplebox/.style={box,draw=nekpurple,fill=neklightpurple},
  graybox/.style={box,draw=nekgray,fill=neklightgray},
  phase/.style={rounded corners=3pt,draw=nekgray,line width=0.8pt,
    fill=white,align=left,inner sep=6pt,font=\small,text width=4.2cm},
  tinylabel/.style={font=\scriptsize,align=center},
  smalllabel/.style={font=\small,align=center}
}

The Dong high-order outflow boundary condition, \texttt{HOutflowBndCondOp}
of the incompressible Navier--Stokes solver, is a coupled condition: it
computes a pressure \emph{Dirichlet} value from an explicitly extrapolated
velocity, solves the pressure Poisson problem, and then computes a velocity
\emph{Neumann} value from that solved pressure. A backflow-selective energy
term stabilises vortical structures as they cross an open boundary. This
chapter follows one time-integration stage through mathematical evaluation,
volume-to-boundary data movement, coefficient projection, and the two global
linear solves. The operator is implemented in
\path{HOutflowBndCondOpImpl.hpp}, in the directory
\path{BoundaryConditions/HOutflowBndCond} of the solver, for the Serial, AVX
and Device execution spaces.

\begin{center}
\resizebox{\linewidth}{!}{%
\begin{tikzpicture}[node distance=5mm]
  \node[bluebox,text width=3.5cm] (a) {extrapolated velocity\\$\vect u^*$ and $\grad\vect u^*$};
  \node[purplebox,right=of a,text width=3.5cm] (b) {pressure boundary\\$p|_{\Gammao}$};
  \node[greenbox,right=of b,text width=3.5cm] (c) {solved pressure\\$p^{n+1}$};
  \node[orangebox,right=of c,text width=3.5cm] (d) {velocity boundary\\$\partial_n\vect u^{n+1}$};
  \draw[dataflow] (a)--(b);
  \draw[stateflow] (b)--(c);
  \draw[dataflow] (c)--(d);
\end{tikzpicture}}
\end{center}

\section{What problem does \HOutflow{} solve?}

An ordinary outlet often prescribes a pressure and assumes zero normal
velocity gradient.  That combination can reflect or destabilise vortices
when part of the outlet experiences local backflow.  Dong's condition instead
couples pressure and velocity through an energy flux that activates smoothly
when the normal velocity changes sign \cite{DongKC2014}.

Let $\vect n$ be the outward unit normal, $\nu>0$ the kinematic viscosity,
$\vect u^*$ the explicitly extrapolated boundary velocity, and $\vect f_b$
the vector expression supplied in the velocity boundary-condition
\code{VALUE}.  On an outflow boundary $\Gammao$, the redesign implements
\begin{align}
  p
  &= \nu\,\vect n\cdot(\grad\vect u^*)\vect n
     - \bigl(\vect E(\vect n,\vect u^*)-\vect f_b\bigr)\cdot\vect n,
     \label{houtflow:eq:pressure}\tag{P}\\
  \frac{\partial\vect u^{n+1}}{\partial n}
  &= \frac{1}{\nu}
     \left\{
       \bigl[p-\nu\,\grad\cdot\vect u^*\bigr]\vect n
       +\vect E(\vect n,\vect u^*)-\vect f_b
     \right\}.
     \label{houtflow:eq:velocity}\tag{V}
\end{align}
Equation~\eqref{houtflow:eq:pressure} is evaluated \emph{before} the pressure solve;
equation~\eqref{houtflow:eq:velocity} is evaluated \emph{after} it.  The divergence
term in \eqref{houtflow:eq:velocity} is essential: omitting it can lock the velocity
at the open boundary.

\subsection{Backflow switch and generalised energy flux}

Define the normal velocity and smooth switch
\begin{equation}
  u_n=\vect n\cdot\vect u^*,
  \qquad
  S_0(u_n)=\frac12\left[1-\tanh\left(\frac{u_n}{U_0\delta}\right)\right].
  \label{houtflow:eq:switch}
\end{equation}
The energy vector is
\begin{equation}
  \boxed{
  \vect E(\vect n,\vect u^*)=
  \left[
    \frac{\theta+\alpha_2}{2}\lvert\vect u^*\rvert^2\vect n
    +\frac{1-\theta+\alpha_1}{2}u_n\vect u^*
  \right]S_0(u_n).}
  \label{houtflow:eq:energy}
\end{equation}

\begin{figure}[htbp]
\centering
\begin{tikzpicture}[>=Latex,x=1cm,y=1cm]
  \draw[->,nekgray] (-4.5,0)--(4.7,0) node[right,font=\small] {$u_n/(U_0\delta)$};
  \draw[->,nekgray] (0,-0.25)--(0,3.3) node[above,font=\small] {$S_0$};
  \draw[nekgray!35,dashed] (-4.3,3)--(4.3,3);
  \draw[nekgray!35,dashed] (-4.3,1.5)--(4.3,1.5);
  \node[left,font=\scriptsize] at (0,3) {$1$};
  \node[left,font=\scriptsize] at (0,1.5) {$1/2$};
  \draw[very thick,nekred,smooth,domain=-4.2:4.2,samples=100]
    plot (\x,{1.5*(1-tanh(\x))});
  \node[redbox,text width=3.0cm] at (-3.0,2.2)
    {local backflow\\$u_n<0$: energy term on};
  \node[bluebox,text width=3.0cm] at (3.0,0.8)
    {outgoing flow\\$u_n>0$: energy term off};
  \draw[flow] (-2.0,2.2)--(-1.2,2.65);
  \draw[flow] (2.0,0.8)--(1.2,0.35);
\end{tikzpicture}
\caption{The switch is continuous, strongly activating the stabilising
energy flux under backflow and suppressing it for outward flow.  Its
transition width is $U_0\delta$.}
\label{houtflow:fig:switch}
\end{figure}

\section{Physical picture at an open boundary}

The boundary condition is evaluated pointwise.  A single outlet can therefore
contain outward-flow and backflow patches at the same time.  The switch in
Figure~\ref{houtflow:fig:switch} responds to the local value of $u_n$, not to a
region-averaged flow direction.

\begin{figure}[htbp]
\centering
\begin{tikzpicture}[>=Latex,scale=0.95]
  \fill[neklightblue] (0,0) rectangle (11,5.2);
  \draw[line width=0.9pt,nekblue] (0,0) rectangle (11,5.2);
  \draw[line width=4pt,nekred] (11,0)--(11,5.2);
  \node[rotate=90,nekred,font=\bfseries\small] at (11.45,2.6)
    {$\Gammao$: \HOutflow{}};

  \foreach \y in {0.65,1.35,2.05,3.85,4.55}
    \draw[-{Latex[length=2.4mm]},nekblue,line width=0.9pt]
      (7.4,\y)--(10.65,\y);
  \foreach \y in {2.55,3.15}
    \draw[-{Latex[length=2.4mm]},nekred,line width=1.05pt]
      (10.65,\y)--(8.75,\y);

  \draw[nekpurple,line width=1.15pt,
    decoration={snake,amplitude=2.2mm,segment length=8mm},decorate]
    (1.0,2.8)--(8.4,2.8);
  \draw[->,nekgreen,line width=1.1pt] (11,4.55)--(12.2,4.55)
    node[right,font=\small] {$\vect n$};

  \node[bluebox,text width=3.0cm] at (4.0,4.25)
    {outgoing patches\\$u_n>0,\ S_0\simeq0$};
  \node[redbox,text width=3.0cm] at (6.0,1.15)
    {vortex-induced backflow\\$u_n<0,\ S_0\simeq1$};
  \draw[flow] (5.2,3.8)--(8.4,3.85);
  \draw[flow] (7.3,1.45)--(9.6,2.55);
\end{tikzpicture}
\caption{A vortex crossing the outlet can produce a local reversal even
when the net flow is outward.  The energy correction acts only where the
normal component is negative.}
\label{houtflow:fig:physical}
\end{figure}

\section{One time-integration stage: the complete workflow}

The public operator, \code{HOutflowBndCondOp}, is not itself a boundary
condition operator.  It holds the boundary data both conditions share and
applies them with two methods: \code{ApplyPressure} sets the Dirichlet pressure
coefficients, and \code{ApplyVelocity} adds the velocity flux.  Two adapters
present these to the linear systems, which accept only Dirichlet and Neumann
condition operators: \code{GetPressureDirichletOp} returns a
\code{DirBndCondOp}, \code{HOutflowPressureAdapter}, added to the pressure
system with \code{AddDirichletBCOp}, and \code{GetVelocityNeumannOp} a
\code{NeuBndCondOp}, \code{HOutflowVelocityAdapter}, added to the velocity
system with \code{AddNeumannBCOp}.  The \code{Apply} of each forwards to the
matching method of the operator, which the operator hands each adapter with
\code{SetOwner} when it creates it, so both refer to the same geometry,
normals, and intermediate arrays.  Figure~\ref{houtflow:fig:stage} shows the
exact ordering inside \code{SolveUnsteadyStokesSystem} of the class
\code{VelocityCorrectionScheme}.

\begin{figure}[p]
\centering
\resizebox{\linewidth}{!}{%
\begin{tikzpicture}[node distance=5mm and 10mm]
  \node[bluebox,text width=14.5cm] (hist) {\textbf{1. IMEX extrapolation}\\
    Histories contain integrator-scaled velocity.  Extrapolate to
    $\vect V^*$, then apply $s=1/(\gamma\Delta t)$ so
    $\vect u^*=s\vect V^*$ and $\grad\vect u^*=s\grad\vect V^*$.};

  \node[bluebox,below=of hist,text width=14.5cm] (gather) {\textbf{2. Restricted volume work}\\
    Gather only velocity elements touching $\Gammao$; evaluate
    $\grad\vect V^*$ with \code{PhysDerivOp}; trace the velocity from the
    full field and all gradient components from those elements to boundary
    quadrature points.};

  \node[purplebox,below=of gather,text width=14.5cm] (pre) {\textbf{3. Pressure boundary update}\\
    Evaluate $\vect f_b(t)$, $u_n$, $S_0$, $\vect E$, $\grad\cdot\vect u^*$,
    and $\vect n\cdot(\grad\vect u^*)\vect n$.  Form (P), forward-transform
    the boundary values, and set the pressure Dirichlet coefficients.};

  \node[greenbox,below=of pre,text width=14.5cm] (psolve) {\textbf{4. Pressure Poisson solve}\\
    The standard Dirichlet adapter first evaluates ordinary XML data;
    \HOutflow{} then overwrites only its tagged pressure coefficients.
    Solve for $p^{n+1}$ and backward-transform it to physical space.};

  \node[orangebox,below=of psolve,text width=14.5cm] (vel) {\textbf{5. Velocity boundary update}\\
    \code{SetPressure} hands the operator the solved pressure field.  When
    the velocity solve applies the paired Neumann adapter, the adapter traces
    $p^{n+1}$ to $\Gammao$ and combines it with the retained
    $\vect E-\vect f_b$ and $\grad\cdot\vect u^*$ to form (V), projects it
    with \code{IProductWRTBaseOp} and scatter-adds it.};

  \node[redbox,below=of vel,text width=14.5cm] (vsolve) {\textbf{6. Velocity solve}\\
    Apply the standard XML Neumann load on every other region, and the
    paired \HOutflow{} load of step 5 on $\Gammao$.  Solve either Helmholtz
    (semi-implicit) or linear ADR (linear-implicit) for $\vect u^{n+1}$.};

  \foreach \a/\b in {hist/gather,gather/pre,pre/psolve,psolve/vel,vel/vsolve}
    \draw[flow] (\a)--(\b);

  \begin{scope}[on background layer]
    \node[fit=(hist)(gather),rounded corners=5pt,fill=nekblue!4,
      draw=nekblue!25,inner sep=7pt] {};
    \node[fit=(pre)(psolve),rounded corners=5pt,fill=nekpurple!4,
      draw=nekpurple!25,inner sep=7pt] {};
    \node[fit=(vel)(vsolve),rounded corners=5pt,fill=nekorange!4,
      draw=nekorange!25,inner sep=7pt] {};
  \end{scope}
\end{tikzpicture}}
\caption{One complete implicit stage.  The retained state created in step 3
is consumed in step 5; recomputing it after the pressure solve is neither
necessary nor desirable.}
\label{houtflow:fig:stage}
\end{figure}

\subsection{Why the stage scale multiplies velocity and its gradient}

The history field passed to \code{UpdatePressure} is stored at the
time-integrator scale.  If it is denoted by $\vect V^*$, the scale set by
\code{SetScale} and the velocity the operator evaluates are
\begin{equation}
  s=\frac{1}{\gamma\Delta t},\qquad
  \vect u^*=s\vect V^*,\qquad
  \grad\vect u^*=s\grad\vect V^*.
\end{equation}
Scaling only the velocity would make the viscous normal derivative and
divergence inconsistent with the switch and energy term.  The boundary
kernel therefore applies \code{velocityScale} to both extracted quantities.

\section{Setup performed once}

Before time stepping, \code{SetVelocityContext} validates the session and
builds reusable geometry and mapping data.  No XML extensions are introduced.

\begin{figure}[htbp]
\centering
\begin{tikzpicture}[node distance=6mm and 10mm]
  \node[graybox,text width=3.3cm] (xml) {read pressure and velocity boundary maps};
  \node[purplebox,right=of xml,text width=3.5cm] (pair) {select pressure\\Dirichlet + \code{HOutflow}};
  \node[orangebox,right=of pair,text width=3.7cm] (valid) {require every velocity component\\Neumann + \code{HOutflow}};
  \node[bluebox,below=11mm of pair,text width=3.8cm] (geom) {match boundary faces to velocity and pressure volume owners by geometry ID};
  \node[greenbox,left=of geom,text width=3.4cm] (restrict) {build restricted list of boundary-adjacent velocity elements};
  \node[graybox,right=of geom,text width=3.8cm] (maps) {build boundary extractions, normals, projection operators, and CG maps};

  \draw[flow] (xml)--(pair);
  \draw[flow] (pair)--(valid);
  \draw[flow] (valid.south) |- (geom.east);
  \draw[flow] (geom)--(restrict);
  \draw[flow] (geom)--(maps);
\end{tikzpicture}
\caption{Setup and validation.  A tag on only one side of the coupled
condition is rejected, as is a tagged pressure Neumann condition.  Pressure
and velocity boundary block shapes and element ordering must agree.}
\label{houtflow:fig:setup}
\end{figure}

The cached objects include:
\begin{itemize}
  \item indexed gathers from the full velocity field to the restricted field;
  \item per region, the extraction of the velocity from the velocity field,
        of its gradient from the restricted field, and of the pressure from
        the pressure field, onto the boundary quadrature points, each an
        \code{MultiRegions::detail::PhysTraceExtractBndCond}
        (\S\ref{sec:phystraceextractbndcond}). The restricted list is built
        from an expansion-information map and has no mesh graph, so the
        gradient extraction is given the mesh graph of the velocity field.
        The pressure and velocity boundary expansions are checked to have the
        same points and block layout, so the extracted fields pair with the
        normals and $\vect f_b$ point by point;
  \item outward normals on each selected boundary region;
  \item \code{PhysDerivOp}, pressure \code{FwdTransBCOp}, and velocity
        \code{IProductWRTBaseOp} instances;
  \item compact boundary-coefficient maps, orientation signs, unique scatter
        targets, and grouped duplicate-target contributions.
\end{itemize}

For MPI partitions, an all-reduce ensures every rank participates in
collective velocity-component assembly-map construction when any rank owns an
\HOutflow{} face.  A rank with no local outlet coefficients then remains a
valid no-op participant.

\section{Pressure phase in detail}

\subsection{Boundary evaluation}

The velocity is extracted onto the boundary quadrature points from the full
velocity field, and its gradient from the restricted elements, in one call
each per region. At each point, the pressure kernel then computes
\begin{align}
  u_n &= n_i u_i^*, &
  \lvert\vect u^*\rvert^2 &= u_i^*u_i^*,\\
  d &= \grad\cdot\vect u^*=\partial_i u_i^*, &
  q &= \vect n\cdot(\grad\vect u^*)\vect n
      =n_i(\partial_j u_i^*)n_j.
\end{align}
With $\vect G=\vect E-\vect f_b$, the state retained for the later velocity
phase is simply
\begin{equation}
  \boxed{\vect G,\quad d},
\end{equation}
while the pressure boundary value is $p=\nu q-\vect G\cdot\vect n$.

\begin{figure}[htbp]
\centering
\resizebox{\linewidth}{!}{%
\begin{tikzpicture}[node distance=5mm and 8mm]
  \node[bluebox,text width=2.7cm] (full) {full velocity field\\$\vect V^*$};
  \node[bluebox,right=of full,text width=2.8cm] (res) {restricted field\\outlet-adjacent elements};
  \node[bluebox,right=of res,text width=2.5cm] (deriv) {volume derivative\\$\grad\vect V^*$};
  \node[purplebox,right=of deriv,text width=2.8cm] (trace) {trace interpolation\\$\vect u^*,\grad\vect u^*$};
  \node[redbox,below=11mm of trace,text width=3.0cm] (kernel) {pointwise kernel\\$S_0,\vect E,\vect G,d,q,p$};
  \node[purplebox,left=of kernel,text width=2.9cm] (fwd) {strong projection\\\code{FwdTransBC}};
  \node[greenbox,left=of fwd,text width=2.9cm] (compact) {compact pressure boundary coefficients};
  \node[greenbox,left=of compact,text width=2.7cm] (set) {CG map/sign\\Dirichlet set};

  \draw[dataflow] (full)--node[above,tinylabel]{indexed gather}(res);
  \draw[dataflow] (res)--(deriv);
  \draw[dataflow] (deriv)--node[above,tinylabel]{gradient}(trace);
  \draw[dataflow] (full.north) to[bend left=12]
    node[above,tinylabel]{velocity}(trace.north);
  \draw[dataflow] (trace)--(kernel);
  \draw[dataflow] (kernel)--(fwd);
  \draw[dataflow] (fwd)--(compact);
  \draw[dataflow] (compact)--(set);

  \node[orangebox,below=8mm of kernel,text width=3.1cm] (fb) {$\vect f_b(t)$ from velocity XML expressions};
  \node[graybox,left=of fb,text width=2.8cm] (norm) {cached normals\\$\vect n$};
  \draw[flow] (fb)--(kernel);
  \draw[flow] (norm)--(kernel);
\end{tikzpicture}}
\caption{Pressure-update pipeline.  Blue operations act on volume storage;
purple/red operations act on boundary quadrature data; green operations
produce the strong pressure Dirichlet lift.}
\label{houtflow:fig:pressure-pipeline}
\end{figure}

\subsection{Strong Dirichlet coefficient path}

The physical pressure values on $\Gammao$ are transformed with
\code{FwdTransBCOp}.  They are gathered into a compact region-major array and
then mapped into the domain coefficient field with the CG orientation signs.
This is a \emph{set} operation: after the standard pressure Dirichlet adapter
has evaluated ordinary XML conditions, the \HOutflow{} pressure adapter
overwrites only its selected coefficients.

\section{Velocity phase in detail}

After the Poisson solve, \code{SetPressure} hands the operator the solved
pressure field, which it does not own.  The velocity flux is not computed
ahead of the velocity solve: the paired Neumann adapter computes it in its
\code{Apply}, which the linear system calls once per solve, from the values
the pressure field holds at that moment.  The flux applied can therefore
never be older than the pressure, and no separate update call has to be
ordered between the pressure and velocity solves; the adapter's
\code{UpdateBndCoeffs} has nothing to do.  Inside \code{Apply}, the solved
pressure is extracted from its boundary-adjacent pressure elements onto the
same boundary quadrature points.
Using retained $\vect G$ and $d$, the pointwise target is
\begin{equation}
  \vect h_{\mathrm{desired}}
  =\frac{1}{\nu}\left[(p-\nu d)\vect n+\vect G\right]
  =\frac{\partial\vect u^{n+1}}{\partial n}.
  \label{houtflow:eq:desired}
\end{equation}

\subsection{Why the default adapter leaves the outlet out}

The velocity values of the \HOutflow{} regions are interpreted as
$\vect f_b$, which enters equations~\eqref{houtflow:eq:pressure}
and~\eqref{houtflow:eq:velocity} through $\vect G=\vect E-\vect f_b$, and
so $\vect h_{\mathrm{desired}}$ already.
The velocity adapter is added to the velocity system with
\code{AddNeumannBCOp(op, "HOutflow")}, so the default Neumann operator of
the system, which evaluates the \code{VALUE} of every other velocity
Neumann region, leaves these regions out, and the paired adapter supplies
the whole load
\begin{equation}
  \vect h_{\mathrm{custom}}=\vect h_{\mathrm{desired}}.
  \label{houtflow:eq:velocity-load}
\end{equation}

\begin{figure}[htbp]
\centering
\resizebox{\linewidth}{!}{%
\begin{tikzpicture}[node distance=6mm and 8mm]
  \node[greenbox,text width=3.0cm] (p) {solved pressure volume field\\$p^{n+1}$};
  \node[purplebox,right=of p,text width=2.9cm] (ptrace) {pressure trace\\$p|_{\Gammao}$};
  \node[orangebox,right=of ptrace,text width=3.2cm] (combine) {combine with cached\\$\vect G,d,\vect n$};
  \node[orangebox,right=of combine,text width=3.1cm] (adjust) {custom physical load\\$\vect h_{\mathrm{desired}}$};

  \node[purplebox,below=11mm of adjust,text width=3.1cm] (iprod) {weak projection\\\code{IProductWRTBase}};
  \node[greenbox,left=of iprod,text width=3.1cm] (coeff) {compact velocity boundary coefficients};
  \node[greenbox,left=of coeff,text width=3.2cm] (scatter) {CG unique/group maps\\Neumann scatter-add};
  \node[graybox,left=of scatter,text width=2.8cm] (rhs) {velocity-system RHS};

  \draw[dataflow] (p)--(ptrace);
  \draw[dataflow] (ptrace)--(combine);
  \draw[dataflow] (combine)--(adjust);
  \draw[dataflow] (adjust)--(iprod);
  \draw[dataflow] (iprod)--(coeff);
  \draw[dataflow] (coeff)--(scatter);
  \draw[dataflow] (scatter)--(rhs);

  \node[bluebox,below=8mm of scatter,text width=3.2cm] (default) {standard Neumann load\\other regions};
  \draw[flow] (default)--(rhs);
\end{tikzpicture}}
\caption{Velocity-update pipeline.  Unlike pressure, velocity boundary data
is a weak load: it is integrated against the boundary basis and scatter-added
to the linear-system right-hand side.}
\label{houtflow:fig:velocity-pipeline}
\end{figure}

Duplicate local CG targets at boundary corners or edges are handled by a
unique-target subset plus grouped contributions.  This avoids both
double-counting and atomics while preserving orientation signs.

\section{Memory and execution-space workflow}

The high-level solver has no backend-specific branches.  Factory registration
selects the Serial, AVX or Device instantiation of
\code{HOutflowBndCondOpImpl}; the
same templated workflow then uses the chosen execution and memory spaces.

\begin{figure}[htbp]
\centering
\resizebox{\linewidth}{!}{%
\begin{tikzpicture}[node distance=7mm]
  \node[graybox,text width=14.7cm] (setup) {\textbf{Host setup, once}\\
    region/tag validation $\rightarrow$ geometry matching $\rightarrow$
    boundary extractions and normals $\rightarrow$ maps copied to execution memory};

  \node[bluebox,below=of setup,text width=14.7cm] (device) {\textbf{Execution memory, every stage}\\
    restricted gather $\rightarrow$ volume derivative $\rightarrow$ trace
    evaluation $\rightarrow$ pointwise kernels $\rightarrow$ projections
    $\rightarrow$ CG set/add};

  \node[purplebox,below=of device,text width=6.8cm,xshift=-3.95cm] (shared) {\textbf{Shared boundary state}\\
    normals, $\vect E-\vect f_b$, $\grad\cdot\vect u^*$,
    $\vect f_b$, pressure/velocity work fields};
  \node[greenbox,right=8mm of shared,text width=6.8cm] (streams) {\textbf{Block streams}\\
    events connect restricted derivatives, boundary batches, projections,
    and mapped coefficient gathers};

  \draw[dataflow] (setup)--node[right,tinylabel]{initial transfer}(device);
  \draw[stateflow] (device)--(shared);
  \draw[stateflow] (device)--(streams);
  \draw[stateflow,dashed] (shared.north) to[bend left=24](device.south);
\end{tikzpicture}}
\caption{Host/device boundary.  Dynamic stage data stays in the active
execution memory; the workflow introduces no per-step host round-trip of
boundary fields.}
\label{houtflow:fig:memory}
\end{figure}

The restricted velocity and its gradient take the interleave width of the
velocity, 1 or that of the execution space, so that the volume derivative
reads and writes them without reshaping. The restricted gather copies through
index maps built on the host for the interleave widths of the source and the
restricted block, the first time the pair is seen, and kept, so no kernel
divides an index, and numbers its threads by restricted point, the element
lane varying fastest, so that its writes are coalesced. Each gather runs on
the stream of its restricted block,
which waits on an event recorded on the stream of its source block, and the
source streams wait on the gathers in turn; the host never waits. Block-local
streams then order the derivative of the velocity. Each extraction, of the
velocity, of its gradient and of the solved pressure, orders itself after
the streams of the blocks it reads, and
the streams of the region blocks the pointwise kernels run on after itself,
with events (see \S\ref{sec:phystraceextractbndcond}). The projected boundary
coefficients are copied into a compact array on the streams of the region
blocks and gathered into the layout of each domain block on the stream of
that block; the region and domain block streams wait on one another through
the operator's own stream before and after the gathers. The final set or
scatter-add indexes each block of the domain coefficient field at the
interleave width of the block, 1 or that of the execution space, through
target indices built on the host the first time a block at that width is
seen, so the field is updated in its own layout and never reshaped.

\FloatBarrier
\section{Session configuration}

The session syntax is that of the IncNavierStokesSolver.  Each selected
region must pair velocity
Neumann conditions and a pressure Dirichlet condition, all tagged with
\code{USERDEFINEDTYPE="HOutflow"}.

\begin{lstlisting}[style=XMLStyle,caption={Minimal two-dimensional HOutflow region.}]
<REGION REF="1">
  <N VAR="u" USERDEFINEDTYPE="HOutflow" VALUE="fb_u(x,y,t)" />
  <N VAR="v" USERDEFINEDTYPE="HOutflow" VALUE="fb_v(x,y,t)" />
  <D VAR="p" USERDEFINEDTYPE="HOutflow" VALUE="0" />
</REGION>
\end{lstlisting}

The pressure \code{VALUE} is a placeholder evaluated by the ordinary
Dirichlet adapter and then replaced by equation~\eqref{houtflow:eq:pressure}.  The
velocity values are meaningful: they define the components of $\vect f_b$.

\begin{table}[htbp]
\centering
\caption{Optional session parameters and preserved defaults.}
\begin{tabular}{@{}lll p{6.8cm}@{}}
\toprule
Parameter & Symbol & Default & Role \\
\midrule
\code{OutflowBC\_theta}  & $\theta$  & $1$    & weights the normal kinetic-energy term \\
\code{OutflowBC\_alpha1} & $\alpha_1$& $0$    & weights the $(\vect n\cdot\vect u)\vect u$ term \\
\code{OutflowBC\_alpha2} & $\alpha_2$& $0$    & augments the normal kinetic-energy term \\
\code{U0\_HighOrderBC}   & $U_0$     & $1$    & reference velocity in the switch \\
\code{Delta\_HighOrderBC}& $\delta$  & $1/20$ & non-dimensional switch transition width \\
\bottomrule
\end{tabular}
\end{table}

$U_0$ and $\delta$ must be positive.  In three dimensions, all three velocity
components must carry paired \HOutflow{} Neumann conditions.

\section{Comparison with ordinary higher-order pressure boundaries}

\begin{table}[htbp]
\centering
\caption{The two operators solve different boundary-condition problems.}
\begin{tabular}{@{}>{\raggedright\arraybackslash}p{3.0cm}
                  >{\raggedright\arraybackslash}p{5.0cm}
                  >{\raggedright\arraybackslash}p{5.4cm}@{}}
\toprule
 & HOPressure & Dong \HOutflow{} \\
\midrule
Pressure condition & Neumann, from momentum consistency & Dirichlet, from
normal viscous stress and energy flux \\
Velocity condition & prescribed Dirichlet & computed Neumann \\
Primary explicit data & RHS, predictor, and extrapolated velocity &
extrapolated velocity $\vect u^*$ \\
After pressure solve & no paired update & solved pressure enters velocity
condition \\
Backflow stabilisation & none & smooth $S_0$-weighted energy term \\
Linear-system action & pressure RHS scatter-add & pressure strong set;
velocity RHS scatter-add \\
Tag selection & pressure Neumann regions with
\code{USERDEFINEDTYPE="H"} & paired
\code{USERDEFINEDTYPE="HOutflow"} conditions \\
\bottomrule
\end{tabular}
\end{table}

The two operators select disjoint regions. \HOutflow{} requires a pressure
Dirichlet condition, and a pressure Neumann condition tagged
\code{USERDEFINEDTYPE="HOutflow"} is rejected at setup; the higher-order
pressure condition of chapter~\ref{chap:hopressure-bndcond} selects only
pressure Neumann regions tagged \code{USERDEFINEDTYPE="H"}.

\FloatBarrier
\section{Implementation map and invariants}

\begin{table}[htbp]
\centering
\caption{Where each workflow stage lives in the redesign.}
\small
\begin{tabular}{@{}>{\raggedright\arraybackslash}p{3.8cm}
                  >{\raggedright\arraybackslash}p{10.0cm}@{}}
\toprule
Responsibility & Location under \path{solvers/IncNavierStokesSolverRedesign} \\
\midrule
Public paired interface &
\path{BoundaryConditions/HOutflowBndCond/HOutflowBndCondOp.hpp} \\
Portable kernels &
\path{BoundaryConditions/HOutflowBndCond/HOutflowBndCondKernels.hpp} \\
Setup, maps and workflow &
\path{BoundaryConditions/HOutflowBndCond/HOutflowBndCondOpImpl.hpp} \\
Factory registration &
\code{ADD\_OPERATOR} in the solver's \path{CMakeLists.txt} \\
Pressure/velocity call order &
\path{EquationSystems/VelocityCorrectionScheme.cpp} \\
Host formula reference &
\path{UnitTests/TestHOutflowBndCond.cpp} \\
Manufactured regressions &
\path{Tests/HOutflowKovasznay2D.tst} and
\path{Tests/HOutflowBackflow2D.tst} \\
\bottomrule
\end{tabular}
\end{table}

When modifying the workflow, preserve these invariants:
\begin{itemize}
  \item \code{UpdatePressure} precedes the Poisson solve, and the velocity
        solve, whose Neumann adapter computes the velocity flux, follows it.
  \item The kinematic viscosity is set with \code{SetKinvis} immediately
        before \code{UpdatePressure} and again, with the pressure set by
        \code{SetPressure}, immediately before the velocity solve; the time
        is set with \code{SetTime} and the scale with \code{SetScale}
        immediately before \code{UpdatePressure}: the operator is shared
        with the pressure solve.
  \item The same scaled $\vect u^*$ is used in $S_0$, $\vect E$, the normal
        derivative, and the divergence.
  \item $\vect E-\vect f_b$ and $\grad\cdot\vect u^*$ survive the pressure
        solve without a host round-trip.
  \item Pressure uses a strong forward transform and Dirichlet set; velocity
        uses a boundary inner product and Neumann scatter-add.
  \item The default Neumann operator of the velocity system leaves the
        \HOutflow{} regions out, and the paired adapter applies the whole
        load of equation~\eqref{houtflow:eq:velocity-load}.
  \item Pressure Neumann regions, with or without the higher-order tag, are
        not touched by this operator.
\end{itemize}

\FloatBarrier
\section{A compact debugging checklist}

\begin{figure}[htbp]
\centering
\begin{tikzpicture}[node distance=6mm and 9mm]
  \node[graybox,text width=3.5cm] (sym) {observed outlet instability or wrong sign};
  \node[purplebox,right=of sym,text width=3.5cm] (tag) {paired tags and BC types correct?};
  \node[bluebox,right=of tag,text width=3.5cm] (scale) {same stage scale on $\vect u$ and $\grad\vect u$?};
  \node[orangebox,below=of tag,text width=3.5cm] (state) {$\vect G$ and divergence retained?};
  \node[greenbox,right=of state,text width=3.5cm] (proj) {strong pressure / weak velocity projection?};
  \draw[flow] (sym)--(tag);
  \draw[flow] (tag)--(scale);
  \draw[flow] (scale.south) |- (state.east);
  \draw[flow] (state)--(proj);
\end{tikzpicture}
\caption{Recommended debugging order.  Most integration mistakes can be
isolated before inspecting the global solver.}
\label{houtflow:fig:debug}
\end{figure}

For a numerical sign check, use a planar outlet with $\vect n=(1,0)$ and a
velocity whose first component changes sign along the boundary.  Then verify
that $S_0\approx0$ on the positive portion, $S_0\approx1$ on the negative
portion, and that the Neumann load the paired adapter assembles equals
equation~\eqref{houtflow:eq:desired}.  The manufactured backflow regression follows
exactly this pattern.

\section{Summary}

\HOutflow{} is best viewed as a two-stage boundary state machine:
\begin{enumerate}
  \item extrapolate $\vect u^*$, evaluate the backflow-aware energy state,
        and impose the resulting pressure Dirichlet value;
  \item solve pressure, combine it with the retained state, and impose the
        resulting velocity Neumann load.
\end{enumerate}
The implementation keeps this state on the active backend, restricts volume
derivatives to outlet-adjacent elements, and reuses the redesign's projection
and CG mapping infrastructure.  The result is one solver-level workflow for
the Serial, AVX and Device execution spaces, with the session syntax of the
IncNavierStokesSolver.
